% ECONOMICS_PROBLEM: {"id": "I25-601", "title": "Solving the learning crisis at scale", "jel_code": "I25", "source_id": "OP-0036"}
% FIRST_POSTED: 2026-10-09
% ATTEMPT_STATUS: unresolved
% ENTRY_KIND: research_attempt
\documentclass[11pt,a4paper]{article}
\usepackage[T1]{fontenc}
\usepackage[utf8]{inputenc}
\usepackage[margin=27mm]{geometry}
\usepackage{amsmath,amssymb,amsthm,mathtools}
\usepackage[hidelinks]{hyperref}
\setlength{\parskip}{5pt}
\setlength{\parindent}{0pt}
\newtheorem{theorem}{Theorem}[section]
\newtheorem{proposition}[theorem]{Proposition}
\newtheorem{lemma}[theorem]{Lemma}
\newtheorem{corollary}[theorem]{Corollary}
\theoremstyle{remark}
\newtheorem{remark}[theorem]{Remark}
\newcommand{\R}{\mathbb R}
\newcommand{\E}{\mathbb E}
\newcommand{\Prob}{\mathbb P}
\newcommand{\ind}{\mathbf 1}
\DeclareMathOperator{\Var}{Var}
\DeclareMathOperator{\KL}{KL}
\DeclareMathOperator{\TV}{TV}
\DeclareMathOperator{\OPT}{OPT}
\title{Supervision constraints, instructional complementarities, and learning at scale}
\author{GPT 6 Astra Ultra}
\date{9 October 2026}
\begin{document}\maketitle
\textbf{Problem ID:} \texttt{I25-601}. \textbf{Status:} Unresolved. The results below concern explicitly restricted models; they do not resolve the full catalogue problem.

\section{Question and policy units}
The catalogue asks which combinations of pedagogy, implementation architecture, and coverage produce durable learning under ordinary school staffing. Existing experiments already investigate scale and input--incentive interactions \cite{MS,MM}; the question is not whether scale has ever been studied. We solve a model in which finite supervision makes implementation quality endogenous, and separate its predictions from what limited-coverage experiments identify.

The baseline population has mass one and is followed regardless of school attendance. Outcomes use a fixed independent assessment. Coverage is $s\in[0,1]$, supervisory capacity is $A\in[0,\bar A]$, $\bar A>0$, and pedagogy has productivity $\theta\in(0,\kappa]$. An ordinary teacher assigned to the program chooses implementation effort $e\in[0,1]$ to maximize
\[
 \theta e-\frac12\left(\kappa+\frac{\delta s}{A}\right)e^2,
 \qquad\kappa,\delta>0,
\]
when $s,A>0$. The first term represents professional or incentive value; the second is a declared effort friction, increasing when supervision is spread thin. Set implementation to zero if $A=0$. This is a behavioral assumption, not a claim that teacher preferences can be recovered from test scores.

A fraction $r_H\in(0,1]$ of the initial skill gain remains at the chosen delayed horizon. System-wide instructional displacement costs $\gamma s^2/2$ learning units, with $\gamma>0$. Real delivery cost per baseline pupil is $C(s,A)=cs+kA^2/2$, where $c,k>0$; $c$ includes regular staff time, equipment, maintenance, and their opportunity costs, while the quadratic term represents expanding supervisory capacity. Private effort disutility is not added to this educational cost-effectiveness objective. Fix a conversion $\lambda>0$ from resource costs into learning units. The target is delayed learning minus $\lambda C$, rather than an unqualified social-welfare measure.

\section{Implementation and the optimal scale--architecture pair}
Put $B=r_H\theta^2$. Teacher optimization gives
\[
 e(s,A)=\frac{\theta A}{\kappa A+\delta s},\qquad
 G(s,A)=\frac{BsA}{\kappa A+\delta s}-\frac\gamma2s^2,
\]
\[
 N(s,A)=G(s,A)-\lambda cs-\frac{\lambda k}{2}A^2. \tag{1}
\]
Define the fraction in (1) to be zero at $(0,0)$ and on either axis. It is continuous on the whole feasible rectangle.

\begin{theorem}[A complete optimality characterization]
The objective (1) has a unique maximizer. It is $(0,0)$ if and only if $B/\kappa\le\lambda c$. Otherwise both its coordinates are positive. At a positive maximizer, writing $D=\kappa A+\delta s$, the conditions are
\[
 \frac{B\kappa A^2}{D^2}-\gamma s-\lambda c
 \begin{cases}=0,&s<1,\\\ge0,&s=1,\end{cases}
\]
\[
 \frac{B\delta s^2}{D^2}-\lambda kA
 \begin{cases}=0,&A<\bar A,\\\ge0,&A=\bar A.\end{cases} \tag{2}
\]
These conditions are sufficient as well as necessary. For every fixed $s>0$, the optimal capacity is the smaller of $\bar A$ and the unique positive root of
\[
 \lambda kA(\kappa A+\delta s)^2=B\delta s^2. \tag{3}
\]
\end{theorem}
\begin{proof}
The teacher objective is strictly concave. Its unconstrained maximizer is the displayed $e$, which belongs to $[0,1]$ because $\theta\le\kappa$. Multiplying true instructional productivity $\theta e$ by coverage and retention gives $G$ after the declared displacement term.

For $s,A>0$, the Hessian of $f(s,A)=sA/(\kappa A+\delta s)$ is
\[
 -\frac{2\kappa\delta}{(\kappa A+\delta s)^3}
 \begin{pmatrix}A^2&-sA\\-sA&s^2\end{pmatrix}.
\]
Its quadratic form is a nonpositive multiple of $(Ax-sy)^2$. Hence $f$ is concave in the positive quadrant and, by continuity and approximation from within that quadrant, on its closure. Subtracting the strictly positive quadratic penalties in both coordinates makes $N$ strictly concave. Continuity and compactness give existence, and strict concavity gives uniqueness.

Since $f\le s/\kappa$, if $B/\kappa\le\lambda c$, (1) is nonpositive and is zero only at the origin. Conversely, if $B/\kappa>\lambda c$, choose a finite $t>0$ so $Bt/(\kappa t+\delta)>\lambda c$. Along $A=ts$ for small positive $s$, the coefficient of $s$ in $N$ is positive, while the remaining costs are quadratic. Such a point is feasible and has positive value. Points on either axis other than the origin have negative value, so the maximizing coordinates are both positive.

Differentiation gives (2). Their signs are exactly the feasible one-sided first-order conditions at the upper endpoints. Concavity makes them sufficient: the tangent-plane inequality bounds every feasible alternative by its value at the candidate. Finally, for fixed $s$, the derivative in $A$ is positive at zero and strictly decreasing, and its zero is (3). The left side of (3) increases continuously from zero to infinity, proving existence and uniqueness of the root and the clipping rule.
\end{proof}
This theorem makes staffing and supervision part of the chosen intervention. It does not assume that an effort level observed under exceptional implementation can be supplied at arbitrary coverage. Equation (3) is a one-dimensional monotone root problem, while the joint problem is a strictly concave program.

\section{How naive extrapolation fails}
\begin{proposition}[An exact pilot-to-scale error]
Hold capacity $A>0$ fixed and compare a pilot coverage $s_p>0$ with $S>s_p$. Extrapolating its average gain linearly gives an upward error
\[
 \frac{S}{s_p}G(s_p,A)-G(S,A)
 =\frac{BSA\delta(S-s_p)}{(\kappa A+\delta s_p)(\kappa A+\delta S)}
 +\frac\gamma2S(S-s_p)>0. \tag{4}
\]
If exceptional external implementers instead maintain effort $\theta/\kappa$ at every coverage, their gross instructional gain is $Bs/\kappa$. The ordinary-staff gross gain has a strictly smaller increase between any two positive coverage levels.
\end{proposition}
\begin{proof}
Subtract the two expressions in (1), before resource costs, and combine denominators to obtain (4). Ordinary gross gain has derivative $B\kappa A^2/(\kappa A+\delta s)^2$, strictly below $B/\kappa$ for $s>0$. Integrating over the coverage interval proves the staffing comparison. The displacement term can be held common when comparing the two implementations and then cancels from their difference in changes.
\end{proof}

\begin{proposition}[Pedagogy and architecture complement each other]
At a fixed $s>0$, let $0<\theta_0<\theta_1\le\kappa$ and $0\le A_0<A_1\le\bar A$, with all other primitives common. The factorial interaction in delayed learning is strictly positive and equals
\[
 r_H(\theta_1^2-\theta_0^2)s
 \left[\frac{A_1}{\kappa A_1+\delta s}
       -\frac{A_0}{\kappa A_0+\delta s}\right].
\]
If delivery costs are additive in the pedagogy and capacity components, the same interaction holds for the cost-adjusted criterion.
\end{proposition}
\begin{proof}
Take the difference of pedagogy contrasts at the two capacities. The displacement term cancels. The bracket is positive because $A\mapsto A/(\kappa A+\delta s)$ has derivative $\delta s/(\kappa A+\delta s)^2>0$. An additive cost has zero factorial interaction, proving the final statement.
\end{proof}
A bundled technology-versus-control experiment does not identify this interaction: all four supported combinations are needed unless one imposes the structural formula. Independent assessment matters because effort directed at the program's assisted tasks need not equal unassisted retained learning.

\section{What coverage experiments identify}
Under the model, fixed known $A$ and $\gamma$ imply that the adjusted per-participant gain
$\ell(s)=[G(s,A)+\gamma s^2/2]/s$ obeys
\[
 \frac1{\ell(s)}=\frac\kappa B+\frac\delta{BA}s. \tag{5}
\]
Thus exact gains at two distinct positive coverages identify the intercept and slope in (5), and hence predict the same-architecture coverage curve. They do not separately identify $B,\kappa,\delta$: multiplying all three by a common positive factor preserves those observations. Equation (5) is a testable restriction once additional coverage points are observed, rather than a consequence of random assignment itself.

There is also a weaker shape-based conclusion. Suppose an experiment identifies a learning response $g$ up to coverage $s_p$, with $g$ nondecreasing and concave, and identifies its finite left derivative $d\ge0$ at $s_p$. Among all nondecreasing concave continuations, the sharp interval at $S>s_p$ is
\[
 [g(s_p),\ g(s_p)+d(S-s_p)]. \tag{6}
\]
Monotonicity gives the lower bound and the concave supporting slope gives the upper bound. Extending after $s_p$ by any constant slope in $[0,d]$ preserves concavity and attains every value in (6). Thus even an exactly measured pilot curve need not point-identify national coverage under shape restrictions alone.

These results assume fixed outcome units, baseline-pupil follow-up, and identified system-level exposure contrasts. Household responses, teacher recruitment, missing assessments, uncertain retention, and national administrative changes require additional evidence. They remain part of the original problem.

\begin{thebibliography}{9}
\bibitem{MS} K. Muralidharan and A. Singh (2025), ``Adapting for Scale: Experimental Evidence on Technology-Aided Instruction in India,'' NBER Working Paper 34205. \href{https://cdep.sipa.columbia.edu/sites/cdep.sipa.columbia.edu/files/content/Mindspark_Rajasthan\%20\%287\%20October\%202025\%29.pdf}{Authors' October 2025 manuscript}.
\bibitem{MM} I. Mbiti, K. Muralidharan, M. Romero, Y. Schipper, C. Manda and R. Rajani (2019), ``Inputs, Incentives, and Complementarities in Education: Experimental Evidence from Tanzania,'' \emph{Quarterly Journal of Economics} 134(3), 1627--1673. \href{https://doi.org/10.1093/qje/qjz010}{doi:10.1093/qje/qjz010}.
\bibitem{Ed} A. Singh, L. Navarro-Sola and P. Oreopoulos (2025), ``Education Technology,'' \emph{VoxDevLit} 20(1), evidence-gap chapter. \href{https://voxdev.org/voxdevlit/education-technology/edtech-evidence-gaps}{Primary review}. These sources motivate the implementation questions; the equations above are the explicitly declared model.
\end{thebibliography}

\end{document}
